\uuid{CtYn}
\exo7id{4594}
\titre{exo7 4594}
\auteur{quercia}
\organisation{exo7}
\datecreate{2010-03-14}
\isIndication{false}
\isCorrection{true}
\chapitre{Série entière}
\sousChapitre{Etude au bord}
\cpge{mpsi-mp2i.procedes-sommatoires-discrets.b, pcsi.series-numeriques.b, ptsi.series-numeriques.b, mpsi-mp2i.procedes-sommatoires-discrets.a, pcsi.series-numeriques.a, ptsi.series-numeriques.a, mp-mpi.series-entieres.a, pc.series-entieres.a, psi.series-entieres.a, pt.series-entieres.a, mp-mpi.series-entieres.c, pc.series-entieres.b, psi.series-entieres.b, pt.series-entieres.b}
\module{Analyse}
\niveau{L2}
\difficulte{}

\contenu{
\texte{
Soit $(a_n)$ une suite de réels strictement positifs.
On suppose que le rayon de convergence de la série entière
$A(x) = \sum_{n=0}^\infty a_nx^n$ est 1 et que la série diverge pour $x = 1$.
}
\begin{enumerate}
    \item \question{Montrer que $A(x) \to +\infty$ lorsque $x\to1^-$.}
\reponse{Fonction croissante. $\lim_{x\to1^-} A(x) \ge \sum_{n=0}^N a_n$.}
    \item \question{Soit $(b_n)$ une suite telle que $b_n \sim a_n$ et
    $B(x) = \sum_{n=0}^\infty b_nx^n$.
    Montrer que $B(x) \sim A(x)$ pour $x\to1^-$.}
\reponse{Dém de type Césaro.}
\end{enumerate}
}
